\section{Local replacement and the number of zero strata}
\label{sec:surgery}

We now choose the separating filtration so that its zero-dimensional set
has small expected cardinality in each original simplex.  Throughout this
section the graph, complex, parameter space, dimension, and radii are fixed
as in Sections~\ref{sec:metric} and~\ref{sec:filtrations}.  In particular,
the strict inequality~\eqref{eq:strict-volume} holds.  Only the filtration
will depend on a positive number $\epsilon$.

\begin{proposition}\label{prop:zero-strata}
For every $\epsilon>0$ there is an admissible equivariant locally constant
filtration $(Z_j,T_j)_{j=0}^n$ such that, for every original top simplex
$\sigma$ of $K$,
\begin{equation}\label{eq:zero-strata}
 \int_X\#\bigl(Z_0(x)\cap\sigma\bigr)\,d\mu(x)
 \le C_\sigma\epsilon.
\end{equation}
Here $C_\sigma$ depends on the fixed geometric data but is independent of
$\epsilon$ and of the complexity of the chosen triangulations.
\end{proposition}

\subsection{Successive near minimizers}

The comparison with a cut through a ball follows the separating-set
method of Papasoglu \cite[Section 2, especially Lemma 2.5]{Papasoglu2020}
and its iterated and equivariant forms in Alpert
\cite[Lemmas 6--7 and 11--12]{Alpert2025}. We prove the comparisons for our
simplex-dependent areas, retaining the error from each near minimizer.

\begin{proof}[Proof of Proposition~\ref{prop:zero-strata}]
For an equivariant family $Z$ of dimension at most $j$, define
\begin{equation}\label{eq:normalized-area}
 \mathcal I_j(Z)=
 \sum_{\sigma\ \mathrm{top}\bmod H}\frac{1}{|H_\sigma|}
 \int_X a_j^\sigma\bigl(Z(x)\cap\sigma^\circ\bigr)\,d\mu(x).
\end{equation}
The sum runs over representatives of the original top simplices, and
$H_\sigma$ is the finite setwise stabilizer of $\sigma$.  This quantity is
finite for an admissible family.  There are finitely many representatives,
and, on each, local constancy and compactness of $X$ give finitely many
finite polyhedral configurations, each of finite area.

Start with $(Z_n,T_n)=(|K|,K)$.  At step $d=n,\ldots,1$, keep all previously
chosen data fixed and choose the admissible pair $(Z_{d-1},T_{d-1})$ so that
$\mathcal I_{d-1}(Z_{d-1})$ is within $\epsilon$ of the infimum over the
admissible choices at this step.  Lemma~\ref{lem:filtration-exists} ensures
that this class is nonempty.  Its infimum is a finite nonnegative number,
so such a choice exists.  Consequently, for every admissible replacement
$Z'$ at this step,
\begin{equation}\label{eq:near-minimality}
 \mathcal I_{d-1}(Z_{d-1})
 \le \mathcal I_{d-1}(Z')+\epsilon.
\end{equation}
Neither an attained minimum nor a limit of filtrations is used.

\subsection{Disjoint replacements in the parameter and spatial variables}

Fix a center $c\in|K|$.  The set
\[
 \mathcal F_c=
 \{h\in H\setminus\{1\}:L_R(hc)\cap L_R(c)\ne\varnothing\}
\]
is finite.  Indeed such an intersection implies $\rho(hc,c)\le2R$;
the support-distance bound, the finite graph balls and freeness on vertices
then leave only finitely many possibilities for $h$.

There is a finite clopen partition of $X$ into sets $A$ satisfying
$A\cap hA=\varnothing$ for every $h\in\mathcal F_c$.  To construct one,
use the free action and the zero-dimensional Hausdorff topology to choose,
around each $x\in X$, a clopen neighborhood disjoint from its translates
by this finite set.  A finite subcover exists by compactness; the atoms of
its common refinement have the required property.  Thus
\begin{equation}\label{eq:product-disjointness}
 \{hA\times L_R(hc):h\in H\}
 \quad\text{is a pairwise disjoint family.}
\end{equation}
This partition is fixed independently of $\epsilon$.  The assertion
concerns the product of parameter and spatial variables; the sets $hA$
alone need not be disjoint.

Fix one atom $A$ and any clopen subset $B\subseteq A$.  For $r\le s\le R$
put
\begin{equation}\label{eq:local-areas}
 M_j(s)=\int_B\sum_{\sigma\ \mathrm{top}}
 a_j^\sigma\bigl(Z_j(x)\cap\sigma^\circ\cap L_s(c)\bigr)\,d\mu(x),
 \qquad 0\le j\le n.
\end{equation}
For a generic $t\in(r,R)$, also put
\[
 N_{d-1}(t)=\int_B\sum_{\sigma\ \mathrm{top}}
 a_{d-1}^\sigma
 \bigl(Z_d(x)\cap\sigma^\circ\cap S_t(c)\bigr)\,d\mu(x).
\]
These sums involve only finitely many original simplices.  Unlike
\eqref{eq:normalized-area}, they run over all such simplices and carry no
stabilizer weights.

At stage $d$, write $P=Z_d$ and $Z=Z_{d-1}$.  For each $x$, make the
simultaneous replacement
\begin{equation}\label{eq:local-replacement}
 Z'(x)=
 \left(Z(x)\setminus\bigcup_{h:x\in hB}U_t(hc)\right)
 \cup
 \left(P(x)\cap\bigcup_{h:x\in hB}S_t(hc)\right).
\end{equation}
The active closed balls are pairwise disjoint by
\eqref{eq:product-disjointness}.  The families of balls and levels are
locally finite, and their use is conditioned on clopen sets, so the
construction is equivariant and locally constant on every original
simplex.  Apply Lemma~\ref{lem:refinement} with $T=T_d$,
$S=T_{d-1}$, $P=Z_d$, $Z=Z_{d-1}$, and the centers and activation
sets $(hc,hB)$.  Outside a countable set of levels, it represents
$Z'$ by an admissible common refinement: the inserted levels have
the required dimension drop on faces of $T_d$ and of the old finer
triangulation $T_{d-1}$, and the lemma preserves the diameter bound
on components of $P(x)\setminus Z'(x)$.

We next compute the change in~\eqref{eq:normalized-area}.  Genericity
against the faces of $T_{d-1}$ implies that the old set $Z$ meets any
inserted level in dimension at most $d-2$.  When $d=1$, that intersection
is empty.  Hence the removed $(d-1)$-area is unchanged if the open ball is
replaced by its closure, and the inserted cut has zero $(d-1)$-area
overlap with the retained old set.

The normalization in~\eqref{eq:normalized-area} makes the removed area
exactly $M_{d-1}(t)$ and the inserted area exactly $N_{d-1}(t)$.  Here is
the orbit-counting identity behind this assertion.  Let $Q$ denote the
relevant equivariant set, and let $W(c)$ denote either the ball or the
level.  Covariance of area and invariance of $\mu$ give
\begin{align*}
 &\sum_{\sigma\ \mathrm{top}\bmod H}\frac1{|H_\sigma|}
   \sum_{h\in H}\int_{hB}
   a_j^\sigma\bigl(Q(x)\cap\sigma^\circ\cap W(hc)\bigr)\,d\mu(x)
 \\
 &\quad=\sum_{\sigma\ \mathrm{top}\bmod H}\frac1{|H_\sigma|}
   \sum_{h\in H}\int_B
   a_j^{h^{-1}\sigma}
   \bigl(Q(x)\cap(h^{-1}\sigma)^\circ\cap W(c)\bigr)\,d\mu(x)
 \\
 &\quad=\int_B\sum_{\tau\ \mathrm{top}}
   a_j^\tau\bigl(Q(x)\cap\tau^\circ\cap W(c)\bigr)\,d\mu(x).
\end{align*}
The map $h\mapsto h^{-1}\sigma$ takes each value in the simplex orbit
exactly $|H_\sigma|$ times.  All local contributions are nonnegative,
and only finitely many translates meet any fixed representative simplex.
The disjointness and zero-overlap statements above therefore yield
\[
 \mathcal I_{d-1}(Z')-\mathcal I_{d-1}(Z)
 =N_{d-1}(t)-M_{d-1}(t).
\]
Combining this identity with~\eqref{eq:near-minimality} gives
\begin{equation}\label{eq:local-comparison}
 M_{d-1}(t)\le N_{d-1}(t)+\epsilon
 \quad\text{for almost every }t\in(r,R).
\end{equation}

\subsection{Slicing and the strict volume margin}

We have compared the area of a separating set inside a ball with that
of a cut through the preceding set.  Slicing now connects these
comparisons across the successive dimensions.

Work first in one original top simplex $\sigma$ and one of its tree
metrics.  Decompose $Z_d(x)\cap\sigma^\circ$ into affine polyhedral
pieces on which $f_c=\rho(c,\cdot)$ is affine.  On a $d$-dimensional
piece $Q$, its tangential gradient has norm $\alpha\le1$ by
Lemma~\ref{lem:tree}.  Euclidean slicing perpendicular to that gradient
gives
\[
 \int_r^b\operatorname{vol}_{d-1}(Q\cap\{f_c=t\})\,dt
 =\alpha\operatorname{vol}_d(Q\cap\{r<f_c<b\})
 \le\operatorname{vol}_d(Q\cap L_b(c)).
\]
A constant branch contributes at only one level.  Pieces of dimension
less than $d$, including interfaces of this decomposition, have zero
$(d-1)$-slice area outside finitely many exceptional levels.  Averaging
over the tree metrics, summing over original top interiors, and
integrating over $B$ yields
\begin{equation}\label{eq:coarea}
 \int_r^bN_{d-1}(t)\,dt\le M_d(b),\qquad r\le b\le R.
\end{equation}
These operations are legitimate by local finiteness and local constancy;
on the fixed ball only finitely many configurations are involved.
Every comparison has taken place in a single simplex metric before
averaging, so no compatibility of metrics across original faces is needed.
Indeed~\eqref{eq:iterated-drop} bounds $Z_d$ on a proper original face
by dimension $d-1$, and its generic level slice by dimension $d-2$.
Such faces contribute no area to the preceding comparison.  The same
dimension estimate says that $Z_0$ misses every proper original face.

From~\eqref{eq:local-comparison} and~\eqref{eq:coarea},
\[
 M_d(b)\ge\int_r^bM_{d-1}(t)\,dt-\epsilon(b-r).
\]
Induction, beginning with the monotonicity of $M_0$, gives
\[
 M_d(b)\ge
 \frac{(b-r)^d}{d!}M_0(r)
 -\epsilon\sum_{i=1}^d\frac{(b-r)^i}{i!}.
\]
Set
\[
 L=R-r,\qquad A_n=\frac{L^n}{n!},\qquad
 \mathcal E_n=\sum_{i=1}^n\frac{L^i}{i!}.
\]
Since $Z_n=|K|$, the volume bound~\eqref{eq:volume-bound} implies
\begin{equation}\label{eq:iterated-volume}
 A_nM_0(r)\le V_R\mu(B)+\epsilon \mathcal E_n.
\end{equation}

We may now choose
\[
 B=A\cap\{x\in X:Z_0(x)\cap L_r(c)\ne\varnothing\}.
\]
This set is clopen: finitely many original simplices meet the ball, and
their exact filtration data are constant on a common finite clopen
partition.  Although $B$ uses the completed filtration, its use in an
earlier comparison is legitimate.  Inequality~\eqref{eq:near-minimality}
holds for every admissible competitor with the data down to stage $d$
fixed; it imposes no condition on the later stages.  The same $B$ can
therefore be used in every comparison leading to
\eqref{eq:iterated-volume}.

Each $x\in B$ contributes at least one zero stratum, and each such point
lies in a unique original top interior.  Thus $\mu(B)\le M_0(r)$.
Moreover $M_0(r)$ is the integral of the point count over all of $A$,
because that count is zero on $A\setminus B$.  The strict inequality
$V_R<A_n$ now gives the direct expected-count bound
\begin{equation}\label{eq:ball-zero-strata}
 \int_A\#\bigl(Z_0(x)\cap L_r(c)\bigr)\,d\mu(x)
 =M_0(r)\le\frac{\mathcal E_n}{A_n-V_R}\,\epsilon.
\end{equation}

To pass from the ball estimate to \eqref{eq:zero-strata}, cover a fixed
original top simplex $\sigma$ by finitely many open
$r$-balls with fixed centers $c_1,\ldots,c_q$.  Such a cover exists by
compactness of the simplex.  For each center, use the fixed clopen
partition constructed above, and write $p_i$ for its number of atoms.
Sum~\eqref{eq:ball-zero-strata} over these atoms and then over the balls.
Since the corresponding closed balls also cover $\sigma$, this proves
\eqref{eq:zero-strata} with
\[
 C_\sigma=\frac{\mathcal E_n}{A_n-V_R}\sum_{i=1}^q p_i.
\]
All quantities in this constant were fixed independently of $\epsilon$.
In particular, no bound on the number of subdivision simplices or on
the cardinality of an individual zero-dimensional fiber is needed.
\end{proof}
